\section{Możliwości dalszych badań}

Zidentyfikowane w pracy ograniczenia oraz otrzymane rezultaty pozwalają wyznaczyć kierunki dalszych badań. Przede wszystkim, eksperymenty można rozszerzyć o modele o odmiennej liczbie parametrów, w tym modele o wielkiej skali, a także o większą liczbę modeli reprezentujących różne rodziny. Pozwoliłoby to na szerszą weryfikację, w jakim stopniu zaobserwowane zależności wynikają z konkretnej architektury, liczby parametrów czy jego charakterystyki. Zasadne byłoby również przeprowadzenie dalszych prac w środowisku odzwierciedlającym rzeczywiste warunki produkcyjne.
Dalsze badania mogłoby również obejmować analizę wpływu kwantyzacji na skuteczność ataków oraz metod mitygacji.

Zakres powinien zostać poszerzony o ataki pośrednie. Zwiększyć również należy liczbę badanych metod mitygacji oraz przełamywania zabezpieczeń. Pozwoliłoby to na bardziej kompleksową analizę bezpieczeństwa modeli. Istotnym aspektem jest również uwzględnienie wielojęzycznych wariantów szkodliwych instrukcji. Dalsze badania warto rozbudować również o ataki wieloturowe.

W kontekście metody opartej na sterowaniu aktywacjami jednym z kierunków jest udoskonalenie bramki sterującej odpowiedzialnej za wykrywanie szkodliwej intencji zapytania. Potencjalnym rozwiązaniem rozszerzeniem byłoby także opracowanie dynamicznego dostosowywania siły interwencji na bazie oszacowanego stopnia zagrożenia zapytania wejściowego. Pozwoliłoby to zmniejszyć siłę interwencji w przypadku promptów o niskim poziomie zagrożenia, co redukowałoby degradację jakości generowanej odpowiedzi.

W przypadku metody bazującej na zewnętrznym klasyfikatorze, warto przeanalizować wariant, w którym poszczególne fragmenty odpowiedzi są na bieżąco oceniane w trakcie jej generowania. Takie podejście mogłoby ograniczyć opóźnienie wynikające z oczekiwania na pełną generację tekstu i polepszyć responsywność systemu.